\section{Discussion}

\subsection{Main Findings}

This preregistered experiment provides novel evidence that personalization can deter rather than encourage delegation to algorithms. Participants delegated significantly more to an EV-maximizing algorithm than to a personalized algorithm that incorporated their individual preferences in the first decision (66.9\% vs. 55.8\%, p = 0.0413). However, this difference did not persist in subsequent decisions (p = 0.4413). This counterintuitive finding challenges common assumptions about the benefits of personalization in AI systems and has important implications for the design of algorithmic advisors.

\subsection{Theoretical Implications}

\subsubsection{Algorithm Appreciation vs. Aversion}

Our findings contribute to the ongoing debate about algorithm aversion and appreciation \citep{dietvorst2015, logg2019}. While \citet{logg2019} found that people prefer algorithmic judgment over human judgment in objective tasks, and \citet{dietvorst2015} showed that people avoid algorithms after observing errors, our results suggest that the type of algorithm matters. Participants initially preferred an objective EV algorithm over a personalized algorithm, even though the personalized algorithm should theoretically lead to higher expected utility for them.

This finding aligns with \citet{castelo2019}, who showed that algorithm aversion/appreciation depends on task type (objective vs. subjective). In our study, the EV algorithm may have been perceived as more objective and competent, while the personalized algorithm may have been perceived as more subjective and potentially less reliable.

\subsubsection{Role of Personalization in Trust}

Our results challenge the assumption that personalization always increases trust in AI systems. While \citet{longoni2019} found that personalization increases trust in medical AI, and \citet{jussupow2020} suggested that personalization can reduce aversion when AI accounts for unique situations, our study shows that personalization can actually deter delegation in certain contexts.

One possible explanation is that participants may have perceived the EV algorithm as more competent and reliable because it maximizes expected value, which is an objective criterion. In contrast, the personalized algorithm may have been perceived as less transparent or less trustworthy because it incorporated subjective preferences. This aligns with \citet{hoff2015}, who found that perceived competence is a key factor influencing trust in automation.

\subsubsection{MABA-HABA Framework}

Our findings can be interpreted through the MABA-HABA (Machines Are Better At vs. Humans Are Better At) framework proposed by \citet{glikson2020}. The EV algorithm aligns with the MABA principle, as it excels at objective calculations and expected value maximization. The personalized algorithm, however, blurs the boundary between what machines are better at and what humans are better at, potentially creating confusion about the appropriate division of labor.

\subsubsection{Trust and Reliance}

Our study distinguishes between trust as an attitude and reliance as behavior \citep{lee2004}. While we did not directly measure trust attitudes, our behavioral measures of reliance show that participants were less willing to rely on the personalized algorithm, at least initially. This aligns with \citet{kohn2021}, who emphasized the importance of measuring both trust attitudes and reliance behaviors.

\subsection{Practical Implications}

\subsubsection{AI System Design}

Our findings have important implications for the design of AI systems, particularly in financial advising. Designers should carefully consider when to personalize recommendations and when to use objective algorithms. In contexts where objective criteria are available and valued, such as expected value maximization in financial decisions, objective algorithms may be more trusted and more likely to be used.

\citet{holzmeister2023} found similar patterns in delegation to financial advisors, suggesting that our findings may generalize to real-world financial decision-making contexts. Designers of robo-advisors should consider whether personalization enhances or detracts from perceived competence and trustworthiness.

\subsubsection{User Education and Communication}

Our results suggest that users may not fully understand the benefits of personalized algorithms. From an economic perspective, the personalized algorithm should be preferred because it incorporates revealed preferences and leads to higher expected utility. However, participants initially chose the objective EV algorithm, suggesting a misunderstanding of the algorithm's benefits.

This highlights the importance of user education and transparent communication about algorithm capabilities and limitations. \citet{leightmann2023} found that explainability affects trust in AI, and \citet{ning2024} showed that transparency affects advice utilization. Designers should clearly communicate how personalized algorithms work and why they may be beneficial for users.

\subsubsection{Policy and Regulation}

Our findings also have implications for policy and regulation of AI-assisted decision-making. Regulators should consider that personalization may not always increase trust or adoption of AI systems. Guidelines for AI-assisted decision-making should emphasize transparency about algorithm type and expected outcomes, and should encourage designers to carefully consider the trade-offs between personalization and objectivity.

\subsection{Limitations}

\subsubsection{Experimental Design}

Our study has several limitations. First, it was conducted online with limited stakes, which may affect the generalizability of our findings. While participants received real monetary payments, the amounts were relatively small, and the decisions were hypothetical in the sense that participants did not have real money at risk.

Second, our study focused on a specific domain (lottery choices) and a specific type of decision (financial risk-taking). The extent to which our findings generalize to other domains (e.g., medical decisions, hiring decisions) is unclear and should be examined in future research.

Third, our sample consisted of UK residents recruited from the Prolific platform, which may not be representative of the broader population. Cultural factors may influence trust in AI and delegation behavior, and our findings may not generalize to other cultural contexts.

\subsubsection{Measurement}

We did not directly measure trust attitudes or perceived competence, relying instead on behavioral measures of reliance. While this aligns with the distinction between trust and reliance \citep{lee2004}, future research should include both attitudinal and behavioral measures to provide a more complete picture.

We also did not measure participants' understanding of the algorithms or their perceptions of algorithm competence. Future research should examine these factors to better understand the mechanisms underlying our findings.

\subsubsection{One-Shot Interaction}

Our study involved a one-shot interaction with the algorithms, with no long-term relationship or repeated interactions. In real-world settings, users may develop trust in algorithms over time through repeated interactions. Future research should examine how delegation behavior changes over longer time horizons and with repeated interactions.

\subsection{Future Research}

\subsubsection{Mechanisms}

Future research should investigate the mechanisms underlying our finding that personalization deters delegation. Why do participants initially prefer objective algorithms over personalized ones? Is it due to perceived competence, transparency, or some other factor? Experimental manipulations of these factors could help identify the causal mechanisms.

\subsubsection{Design Interventions}

Future research should test interventions to increase trust in personalized algorithms. For example, would providing explanations of how the personalized algorithm works increase delegation? Would highlighting the benefits of personalization (e.g., higher expected utility) increase adoption? \citet{shin2021} found that explainability and causability affect trust in human-AI collaboration, suggesting that such interventions may be effective.

\subsubsection{Extensions}

Future research should extend our findings to other domains and contexts. Would similar patterns emerge in medical decision-making, hiring decisions, or other domains where personalization is common? \citet{longoni2019} found that personalization increases trust in medical AI, suggesting that domain-specific factors may moderate the effects of personalization.

Future research should also examine cross-cultural differences in delegation to algorithms. Cultural factors may influence trust in AI and preferences for personalization, and our findings may not generalize to all cultural contexts.

\subsubsection{Longitudinal Studies}

Future research should conduct longitudinal studies of delegation over time. How does delegation behavior change as users gain more experience with algorithms? Do users learn to trust personalized algorithms over time? Longitudinal studies could provide insights into the dynamics of trust and reliance in human-AI collaboration.

\subsection{Conclusion}

This preregistered experiment provides novel evidence that personalization can deter rather than encourage delegation to algorithms. Participants initially delegated more to an objective EV algorithm than to a personalized algorithm, even though the personalized algorithm should theoretically lead to higher expected utility. This counterintuitive finding challenges common assumptions about the benefits of personalization in AI systems and has important implications for the design of algorithmic advisors.

Our results contribute to the literature on algorithm appreciation, trust in AI, and human-AI collaboration. They suggest that in certain contexts, particularly those involving objective calculations, people may prefer algorithms that are perceived as competent and objective rather than those that attempt to personalize recommendations. This has practical implications for the design of AI systems in finance and other domains where delegation decisions are important.

Future research should explore the mechanisms underlying this effect and test interventions to optimize the balance between personalization and trust. As AI continues to play an increasing role in decision-making, understanding when and why people delegate to algorithms will be crucial for designing systems that are both effective and trusted.